\subsection{Initialisierung des Systems}
Die Initialisierung erfolgt einmalig beim Systemstart. 
Dabei werden alle Hardwarekomponenten konfiguriert, die Netzwerkverbindung aufgebaut sowie die aktuelle Zeit über einen \gls{ntp}-Server synchronisiert. 
Zusätzlich werden erste Sensorwerte und Wetterdaten geladen, um direkt nach dem Start eine gültige Anzeige zu gewährleisten. 
Der Ablauf der Initialisierung ist in Abbildung~\ref{fig:setup} dargestellt.

\begin{figure}[H]
\centering
\begin{tikzpicture}[
    node distance=0.8cm,
    startstop/.style={rectangle, rounded corners, minimum width=4cm, minimum height=0.8cm, draw=black},
    process/.style={rectangle, minimum width=4.8cm, minimum height=0.8cm, draw=black},
    arrow/.style={->, thick}
]

\node (start) [startstop] {Start};
\node (s1) [process, below=of start] {Serielle Schnittstelle starten};
\node (s2) [process, below=of s1] {LED-Streifen initialisieren};
\node (s3) [process, below=of s2] {Sensor initialisieren};
\node (s4) [process, below=of s3] {WLAN verbinden};
\node (s5) [process, below=of s4] {NTP konfigurieren};
\node (s6) [process, below=of s5] {Zeit synchronisieren};
\node (s7) [process, below=of s6] {Wetterdaten abrufen};
\node (s8) [process, below=of s7] {Sensorwerte einlesen};
\node (s9) [process, below=of s8] {Erste Anzeige erzeugen};
\node (s10) [process, below=of s9] {HTTP-Server starten};
\node (end) [startstop, below=of s10] {System bereit};

\draw [arrow] (start) -- (s1);
\draw [arrow] (s1) -- (s2);
\draw [arrow] (s2) -- (s3);
\draw [arrow] (s3) -- (s4);
\draw [arrow] (s4) -- (s5);
\draw [arrow] (s5) -- (s6);
\draw [arrow] (s6) -- (s7);
\draw [arrow] (s7) -- (s8);
\draw [arrow] (s8) -- (s9);
\draw [arrow] (s9) -- (s10);
\draw [arrow] (s10) -- (end);

\end{tikzpicture}
\caption{Schematischer Programmablaufplan der Systeminitialisierung.}
\label{fig:setup}
\end{figure}